\subsection{Cambio di variabile nei limiti}
\label{sec:esercizi-limiti-cambio-di-variabile}

\begin{exercise}[Argomento trigonometrico tendente a uno]
\label{ex:limiti-foglio-2311-21}
\difficoltaesercizio{0}
Calcolare il limite
\[
  \lim_{x\to0}\frac{\sin(\cos x)}{\cos x}.
\]
\rimandoteoria{prop:lez20-limiti-trigonometrici}
\end{exercise}
\begin{proof}[Soluzione dell'esercizio]
Posto $t:=\cos x$, si ha $t\to1$ per $x\to0$.
Per continuità,
\[
  \lim_{x\to0}\frac{\sin(\cos x)}{\cos x}
  =\lim_{t\to1}\frac{\sin t}{t}
  =\sin1.
\]
\end{proof}

\begin{exercise}[Seno di un reciproco a infinito]
\label{ex:limiti-foglio-2311-23}
\difficoltaesercizio{0}
Calcolare il limite
\[
  \lim_{x\to+\infty}x\sin\left(\frac1x\right).
\]
\rimandoteoria{thm:lez20-limite-seno}
\end{exercise}
\begin{proof}[Soluzione dell'esercizio]
Posto $t:=1/x$, si ha $t\to0^+$ per $x\to+\infty$.
Per il limite notevole del seno,
\[
  \lim_{x\to+\infty}x\sin\left(\frac1x\right)
  =\lim_{t\to0^+}\frac{\sin t}{t}
  =1.
\]
\end{proof}

\clearpage

\begin{exercise}[Seno di un polinomio vicino a uno]
\label{ex:limiti-foglio-2311-26}
\difficoltaesercizio{0}
Calcolare il limite
\[
  \lim_{x\to1}\frac{\sin(x^2-1)}{x-1}.
\]
\rimandoteoria{thm:lez20-limite-seno}
\end{exercise}
\begin{proof}[Soluzione dell'esercizio]
Posto $t:=x-1$, si ha $t\to0$ per $x\to1$ e
$x^2-1=t(t+2)$. Per il limite notevole del seno,
\[
  \lim_{x\to1}\frac{\sin(x^2-1)}{x-1}
  =\lim_{t\to0}\frac{\sin(t(t+2))}{t(t+2)}(t+2)
  =1\cdot2=2.
\]
\end{proof}
